% 付録F ノートブックガイド
\section*{付録F\quad ノートブックガイド}
\addcontentsline{toc}{section}{付録F\quad ノートブックガイド}

本書のノートブック NB0--NB8 は \texttt{cmbcore} を用い、各章の数値図を読者が
再現・改変できるよう構成されている。本文章との対応は下表のとおり。

\begin{center}
\begin{tabular}{lll}
\hline
ノートブック & 対応章 & 主な図・内容 \\
\hline
NB0 setup\_and\_units & 序章 & SI定数・単位換算・\texttt{Params} \\
NB1 background & 第1章 & $\Omega_i(x)$, $\mathcal H(x)$, $\eta,z_{\rm eq}$ \\
NB2 recombination & 第3章 & $X_e$（Saha/Peebles）, $\tau,\tilde g$, $z_*$ \\
NB3 perturbations\_single\_k & 第6--9章 & $\Theta_\ell,\delta_c,\Phi,\Psi$ の進化 \\
NB4 transfer\_functions & 第10章 & ソース4項, $\Theta_\ell(k)$, $j_\ell$ \\
NB5 power\_spectrum & 第11--12章 & $\mathcal D_\ell$ vs Planck, ピーク検出 \\
NB6 parameter\_dependence & 第13章 & $\tau$, $z_*$, $z_{\rm re}$ の定量依存 \\
NB7 analytic\_vs\_numeric & 第8--9,12章 & $\ell_A,k_D$, ピーク位置予言 \\
NB8 validation\_class & 第14章 & 特徴量検証・収束・CLASS比較 \\
\hline
\end{tabular}
\end{center}

\paragraph{実行環境}
Python 3.11+、\texttt{numpy/scipy/matplotlib}（必須）、\texttt{nbformat/nbconvert}。
\texttt{pip install -e .} 後に \texttt{make notebooks} で
\texttt{jupyter nbconvert --execute} により全冊を完走できる。重い分光計算は
\texttt{figures/cls\_fiducial.npz} と \texttt{figures/param\_study.npz} の
キャッシュを参照するため高速に再現される。乱数種は \texttt{seed=5220} に固定。
\texttt{classy} がある環境では NB6/NB8 の CLASS 比較セルが有効になる。
